% Beamer starter: Gotham theme.
\IfFileExists{.latex-editing-toolkit/beamer-class-options.tex}{\input{.latex-editing-toolkit/beamer-class-options.tex}}{}
\documentclass{beamer}
\providecommand{\ToolkitBeamerTitle}{Presentation Title}
\providecommand{\ToolkitBeamerAuthor}{Author Name}
\providecommand{\ToolkitBeamerInstitute}{Institute}
\providecommand{\ToolkitBeamerDate}{\today}

\usepackage{amsmath}
\usepackage{hyperref}
\makeatletter
\def\input@path{{beamer/gotham/}}
\makeatother
\usetheme{gotham}
\IfFileExists{.latex-editing-toolkit/beamer-settings.tex}{\input{.latex-editing-toolkit/beamer-settings.tex}}{}

\title{\ToolkitBeamerTitle}
\author{\ToolkitBeamerAuthor}
\institute{\ToolkitBeamerInstitute}
\date{\ToolkitBeamerDate}

\begin{document}

\begin{frame}
  \maketitle
\end{frame}

\section{Getting Started}

\begin{frame}{First Slide}
  Replace this content with the main idea of your presentation.
  \begin{block}{A useful block}
    Gotham supports a clean presentation surface for equations, figures, and concise arguments.
  \end{block}
\end{frame}

\end{document}
